% ============================================================
% PyQCU CUDA-C++ MultiGrid 各层算子构造 — 代码与物理结构整理
% 编译：xelatex 两遍；交付闸门：Overfull=0, Float too large=0
% ============================================================
\documentclass[UTF8,aspectratio=169,11pt]{ctexbeamer}

\usepackage{amsmath,amssymb}
\usepackage{booktabs}
\usepackage{tabularx}
\usepackage{array}
\usepackage{tikz}
\usetikzlibrary{arrows.meta,positioning,calc,fit,shapes.geometric}
\usepackage{xcolor}
\usepackage{microtype}
\usepackage{hyperref}

% ===== 色板：颜色只增强层次，不承载唯一信息 =====
\definecolor{primary}{HTML}{17365D}
\definecolor{accent}{HTML}{1F77B4}
\definecolor{evidence}{HTML}{1E8449}
\definecolor{warning}{HTML}{C55A11}
\definecolor{muted}{HTML}{5B6573}
\definecolor{panel}{HTML}{F4F7FA}
\definecolor{linegray}{HTML}{CBD5E1}
\definecolor{good}{HTML}{EAF5EE}
\definecolor{warnbg}{HTML}{FFF4E8}

\setbeamersize{text margin left=0.55cm,text margin right=0.55cm}
\setlength{\emergencystretch}{2em}
\setbeamertemplate{navigation symbols}{}
\setbeamercolor{normal text}{fg=black!86,bg=white}
\setbeamercolor{structure}{fg=primary}
\setbeamercolor{frametitle}{fg=primary,bg=white}
\setbeamercolor{block title}{fg=primary,bg=panel}
\setbeamercolor{block body}{fg=black!86,bg=panel}
\setbeamerfont{frametitle}{size=\large,series=\bfseries}
\setbeamerfont{normal text}{size=\normalsize}
\setbeamertemplate{frametitle}{%
  \vskip0.12cm\insertframetitle\par\vskip0.08cm}
\setbeamertemplate{footline}{%
  \hbox{\begin{beamercolorbox}[wd=\paperwidth,ht=2.3ex,dp=0.9ex,
    leftskip=0.55cm,rightskip=0.55cm]{author in head/foot}%
    \scriptsize\color{muted}\insertshortauthor\hfill
    \insertshorttitle\hfill\insertframenumber/\inserttotalframenumber
  \end{beamercolorbox}}}

\newcolumntype{Y}{>{\raggedright\arraybackslash}X}
\newcolumntype{P}[1]{>{\raggedright\arraybackslash}p{#1}}
\newcommand{\source}[1]{%
  \vspace{-0.08em}\par{\raggedright\scriptsize\color{muted}来源：#1\par}}
\newcommand{\src}[1]{\nolinkurl{#1}}
\newcommand{\verified}{\textcolor{evidence}{确证}}
\newcommand{\inferred}{\textcolor{accent}{推断}}
\newcommand{\unverified}{\textcolor{warning}{未验证}}

\tikzset{
  flow/.style={draw=accent,fill=accent!8,rounded corners=2pt,align=center,
    minimum height=0.70cm,text width=2.05cm,font=\small},
  state/.style={draw=primary,fill=primary!7,rounded corners=2pt,align=center,
    minimum height=0.64cm,text width=2.55cm,font=\small},
  arrow/.style={-{Latex[length=2mm]},draw=primary,semithick},
  relation/.style={-{Latex[length=1.7mm]},draw=muted,thin,dashed},
  labelstyle/.style={font=\scriptsize,inner sep=1.5pt,fill=white,text=muted}
}

\title[PyQCU CUDA-C++ MG]{PyQCU CUDA-C++ MultiGrid 各层算子构造}
\subtitle{从奇子格 Schur 算子到 33 点 Galerkin 粗算子}
\author[代码与物理结构整理]{代码与物理结构整理}
\institute{PyQCU · Lattice QCD}
\date{2026-08-30}

\begin{document}

\begin{frame}[plain]
  \titlepage
\end{frame}

\section{结论先行}

\begin{frame}[t]{核心结论：层级由 Python 构造，CUDA-C++ 负责应用}
  \begin{columns}[T,totalwidth=\textwidth]
    \begin{column}{0.31\textwidth}
      \begin{block}{Level 0}
        \small
        主求解使用奇子格 Schur 算子
        \[
        S=D_{oo}-\kappa^2H_{oe}D_{ee}^{-1}H_{eo}.
        \]
      \end{block}
    \end{column}
    \begin{column}{0.31\textwidth}
      \begin{block}{粗层}
        \small
        每次转移构造局部正交基 $P_l$，并按
        \[
        A_{l+1}=P_l^\dagger A_lP_l
        \]
        投影为 33 点 stencil。
      \end{block}
    \end{column}
    \begin{column}{0.31\textwidth}
      \begin{block}{后端边界}
        \small
        C++ 初始化只解析层几何、分配求解向量；通过
        \src{set_ptrs[30+4*l:34+4*l]} 接收
        \src{lonv/hnn/hdg/sit}。
      \end{block}
    \end{column}
  \end{columns}
  \vspace{0.25em}
  \begin{tabular}{@{}p{0.94\textwidth}@{}}
    \colorbox{good}{\parbox{0.92\textwidth}{\textbf{当前判断（\verified）：}\quad
      严格说，当前“CUDA-C++ 版 MultiGrid”不是在 C++ 内部生成粗算子；层级资产由
      \src{build_schur_levels()} 生成或从 HDF5 缓存加载，\src{libqcu.so} 在 V-cycle 中消费这些资产。}}
  \end{tabular}
  \source{\src{pyqcu/cuda/_multi_gpu.py:63,111}；\src{cpp/cuda/qcu/src/apply_clover_multigrid.cu:11,55}；\src{cpp/cuda/qcu/include/lattice_clover_multigrid.h:2796}。}
\end{frame}

\begin{frame}[t]{总流程：同一条 Galerkin 递归链贯穿全部层级}
  \begin{center}
    \begin{tikzpicture}[node distance=0.24cm and 0.20cm,scale=0.91,transform shape]
      \node[flow] (d) {规范场 $U$\newline Clover $C$};
      \node[flow,right=of d] (s) {Level 0\newline Schur $S$};
      \node[flow,right=of s] (nv) {随机向量\newline 低模富集};
      \node[flow,right=of nv] (p) {局部 QR\newline 得到 $P_l$};
      \node[flow,right=of p] (g) {Galerkin\newline $P_l^\dagger A_lP_l$};
      \node[flow,right=of g] (a) {Level $l+1$\newline 33 点应用};
      \draw[arrow] (d) -- (s);
      \draw[arrow] (s) -- (nv);
      \draw[arrow] (nv) -- (p);
      \draw[arrow] (p) -- (g);
      \draw[arrow] (g) -- (a);
      \draw[relation] (a.south) |- ++(0,-0.48) -| node[below=1.5pt,labelstyle]{作为下一层 $A_l$} (nv.south);
    \end{tikzpicture}
  \end{center}
  \vspace{-0.05em}
  \begin{tabular}{@{}l@{}}
    $A_0=S,\quad A_1=P_0^\dagger S P_0,\quad
    A_2=P_1^\dagger A_1P_1,\quad\cdots$ \\
    每个转移只保存 $\{P_l,\;\text{on-site},\;\text{nearest},\;\text{diagonal}\}$，不保存完整稠密矩阵。
  \end{tabular}
  \vspace{0.12em}
  \begin{block}{为什么递归可行？}
    \small
    $P_l$ 具有 aggregate 局部支撑，且 $A_l$ 只有有限邻接范围；因此投影后的
    $A_{l+1}$ 仍可用固定宽度 stencil 表示，随后重复同一构造流程。
  \end{block}
  \source{\src{pyqcu/cuda/_multi_gpu.py:101-185}；\src{pyqcu/tools/_multigrid.py:738-770}。图为实现逻辑抽象。}
\end{frame}

\section{细层与转移基}

\begin{frame}[t]{Level 0 不是粗化后的 Wilson 算子，而是奇子格 Schur 算子}
  \begin{columns}[T,totalwidth=\textwidth]
    \begin{column}{0.49\textwidth}
      \setlength{\jot}{2pt}
      \begin{align*}
        D&=\begin{pmatrix}
          D_{ee}&-\kappa H_{eo}\\
          -\kappa H_{oe}&D_{oo}
        \end{pmatrix},\\
        A_0\equiv S&=D_{oo}-\kappa^2H_{oe}D_{ee}^{-1}H_{eo}.
      \end{align*}
      \begin{block}{物理图像}
        \small
        消去偶点后，奇点之间通过“奇 $\to$ 偶 $\to$ 奇”的两次 hopping 相连；
        偶点的局部 Clover 逆保留规范场背景和 spin--color 耦合。
      \end{block}
    \end{column}
    \begin{column}{0.45\textwidth}
      \renewcommand{\arraystretch}{0.82}
      \begin{tabular}{@{}P{0.34\linewidth}P{0.52\linewidth}@{}}
        \toprule
        对象 & 当前实现 \\
        \midrule
        输入/输出 & 奇子格向量 \\
        自由度 & $12=4\times3$ \\
        几何 & $[X,Y,Z,T/2]$ \\
        细层应用 & C++ \src{fine_dslash_op} \\
        构造时 matvec & Python \src{op.matvec_parity} 或 \src{CudaSchurOp} \\
        \bottomrule
      \end{tabular}
      \vspace{0.12em}
      \begin{block}{重要边界}
        \footnotesize
        源码中的 full-site 缓冲和最终 full-op 真残差检查不改变当前层级定义；主循环仍在
        Level 0 调用 Schur 算子。
      \end{block}
    \end{column}
  \end{columns}
  \source{\src{lattice_clover_multigrid.h:815,4531,4824}；\src{_multi_gpu.py:497}。}
\end{frame}

\begin{frame}[t]{每条层间转移先构造局部低模基，再定义 $R_l=P_l^\dagger$}
  \begin{columns}[T,totalwidth=\textwidth]
    \begin{column}{0.48\textwidth}
      \begin{tabular}{@{}l@{}}
        $\text{随机场 }v\in\mathbb{C}^{e_l\times X_l\times Y_l\times Z_l\times T_l}$ \\
        $\quad\downarrow\quad\text{松弛 BiCGStab 近似求 }A_l^{-1}(A_lv)$ \\
        $v\leftarrow v-A_l^{-1}(A_lv)$ \\
        $\quad\downarrow\quad\text{归一化并按 aggregate 重排}$ \\
        $[E_{l+1},e_l,X_c,b_x,Y_c,b_y,Z_c,b_z,T_c,b_t]$ \\
        $\quad\downarrow\quad\text{每个 aggregate 做 batched QR}$ \\
        $P_l\text{ 的局部正交列向量}$
      \end{tabular}
      \vspace{0.10em}
      \begin{block}{低模富集的含义}
        \footnotesize
        代码采用宽松相对容差，保留近零空间方向；这不是最终求解精度，而是为粗空间选择代表性基。
      \end{block}
    \end{column}
    \begin{column}{0.45\textwidth}
      \begin{align*}
        (R_lv)_J(c)&=\sum_{x\in\mathcal A(c),a}
          P_{J,a}^*(x)v_a(x),\\
        (P_lc)_a(x)&=\sum_JP_{J,a}(x)c_J(c).
      \end{align*}
      \begin{block}{局部正交性}
        \small
        QR 在每个 aggregate 内令
        \[
        P_l^\dagger P_l=I
        \]
        （各 aggregate 支撑互不重叠），因此 restriction 不需要额外的全局 Gram 矩阵求逆。
      \end{block}
    \end{column}
  \end{columns}
  \source{\src{tools/_multigrid.py:355,69,112}；\src{src/multigrid.cu:13,73}。}
\end{frame}

\section{Galerkin 粗算子}

\begin{frame}[t]{Galerkin 构造是逐个粗基向量探测，而不是显式形成完整矩阵}
  \begin{columns}[T,totalwidth=\textwidth]
    \begin{column}{0.54\textwidth}
      \begin{tabular}{@{}l@{}}
        $\text{for each coarse site }c\text{ and coarse basis }J:$ \\
        $\quad q\leftarrow e_J(c)$ \\
        $\quad f\leftarrow P_lq$ \\
        $\quad d\leftarrow A_lf$ \\
        $\quad z\leftarrow P_l^\dagger d$ \\
        $\quad \text{将 }z\text{ 在 }c\text{ 及其 32 个邻接位置的块写入 stencil}$
      \end{tabular}
      \vspace{0.4em}
      \begin{block}{列与块的对应}
        \footnotesize
        一个 $(c,J)$ 探针给出 $P_l^\dagger A_lP_l$ 的一列；每个输出位置的系数是
        $E_{l+1}\times E_{l+1}$ 复矩阵块。
      \end{block}
    \end{column}
    \begin{column}{0.40\textwidth}
      \begin{equation*}
        A_{l+1}=P_l^\dagger A_lP_l.
      \end{equation*}
      \begin{tabular}{@{}P{0.36\linewidth}P{0.48\linewidth}@{}}
        \toprule
        代码对象 & 数学对象 \\
        \midrule
        \src{lonv} & $P_l$ \\
        \src{restrict} & $P_l^\dagger$ \\
        \src{matvec} & $A_l$ \\
        \src{sit / hnn / hdg} & $A_{l+1}$ 的非零块 \\
        \bottomrule
      \end{tabular}
      \vspace{0.4em}
      \textcolor{accent}{代码注释常写 $P^T$；复数实现实际使用共轭转置 $P^\dagger$。}
    \end{column}
  \end{columns}
  \source{\src{tools/_multigrid.py:691,738}；\src{lattice_clover_multigrid.h:892}。}
\end{frame}

\begin{frame}[t]{Schur 的两次 hopping 使粗层采用 33 点宽 stencil}
  \begin{columns}[T,totalwidth=\textwidth]
    \begin{column}{0.43\textwidth}
      \begin{align*}
        1&\quad\text{onsite},\\
        4\times2&=8\quad\text{同方向正/负邻居},\\
        \binom{4}{2}\times2\times2&=24\quad\text{两方向弯折邻居},\\[2pt]
        N_{\rm stencil}&=1+8+24=33.
      \end{align*}
      \begin{block}{物理图像}
        \small
        $D_{oo}$ 给出本地点；两次 hopping 回到原点、沿同一方向前进，或沿两个不同方向转弯。
      \end{block}
    \end{column}
    \begin{column}{0.53\textwidth}
      \begin{tabular}{@{}P{0.25\linewidth}P{0.63\linewidth}@{}}
        \toprule
        张量 & 语义与形状 \\
        \midrule
        \src{sit} & onsite：$[E,E,X,Y,Z,T]$ \\
        \src{hop_nn} & 正/负邻居：$[2,4,E,E,X,Y,Z,T]$ \\
        \src{hop_diag} & 方向对与符号：$[2,2,6,E,E,X,Y,Z,T]$ \\
        \midrule
        应用式 & $y_j(c)=\mathrm{sit}_{je}(c)x_e(c)$\newline
        $+\mathrm{hop}_{je}(c)x_e(c\pm\hat\mu)$\newline
        $+\mathrm{diag}_{je}(c)x_e(c\pm\hat\mu\pm\hat\nu)$ \\
        \bottomrule
      \end{tabular}
      \vspace{0.35em}
      \textcolor{warning}{\footnotesize 33 是耦合位置数，不是自由度数；每个位置承载一个 $E\times E$ 块。}
    \end{column}
  \end{columns}
  \source{\src{src/multigrid.cu:858,924}；\src{tools/_multigrid.py:738}。}
\end{frame}

\section{运行时边界}

\begin{frame}[t]{C++ 后端的职责：接收四类层级资产并执行宽 stencil}
  \begin{columns}[T,totalwidth=\textwidth]
    \begin{column}{0.49\textwidth}
      \begin{tabular}{@{}P{0.30\linewidth}P{0.56\linewidth}@{}}
        \toprule
        阶段 & C++ 行为 \\
        \midrule
        \src{init} & 读取每层 $E,X,Y,Z,T$，分配求解向量 \\
        \src{set ops} & 保存四个外部指针 \\
        \src{coarse dslash} & 读取当前层 33 点 stencil \\
        MPI 多 rank & 先交换粗层 halo，再启动 CUDA kernel \\
        粗层求解 & BiStabCG / fused coarse solve，算子不变 \\
        \bottomrule
      \end{tabular}
    \end{column}
    \begin{column}{0.47\textwidth}
      \begin{block}{指针协议}
        \footnotesize
        对每个父层 $l$，令 $\texttt{base}=30+4l$：
        \begin{align*}
          \texttt{set\_ptrs[base+0]}&\to\texttt{null\_vecs},\\
          \texttt{set\_ptrs[base+1]}&\to\texttt{hop\_nn},\\
          \texttt{set\_ptrs[base+2]}&\to\texttt{hop\_diag},\\
          \texttt{set\_ptrs[base+3]}&\to\texttt{sit\_packed}.
        \end{align*}
      \end{block}
      \begin{block}{实现分工（\verified）}
        \footnotesize
        Python 侧决定 $P_l$ 和 $A_{l+1}$ 的数值内容；C++ 侧负责 GPU kernel、scratch、halo 和递归求解。
      \end{block}
    \end{column}
  \end{columns}
  \source{见末页源码索引：bridge:11；header:892。}
\end{frame}

\begin{frame}[t]{当前运行路径：缓存命中与构造加速只改变来源，不改变数学算子}
  \begin{center}
    \begin{tikzpicture}[node distance=0.35cm and 0.30cm,scale=0.92,transform shape]
      \node[state] (cfg) {\src{dof_list}, \src{mg_grid}\newline 层级配置};
      \node[flow,right=of cfg] (hit) {检查\newline HDF5 缓存};
      \node[flow,right=of hit] (load) {命中：加载\newline \src{lonv/hnn/hdg/sit}};
      \node[flow,below=of load] (build) {未命中：null\newline QR + probing};
      \node[flow,left=of build] (mv) {matvec\newline Torch 批量或 C++};
      \node[state,left=of mv] (run) {传入 \src{set_ptrs}\newline C++ V-cycle};
      \draw[arrow] (cfg) -- (hit);
      \draw[arrow] (hit) -- node[above=1.5pt,labelstyle]{是} (load);
      \draw[relation] (hit.south) |- node[right=1.5pt,labelstyle]{否} (build.north);
      \draw[arrow] (build) -- (mv);
      \draw[arrow] (load.south) |- (run.north);
      \draw[arrow] (mv) -- (run);
    \end{tikzpicture}
  \end{center}
  \vspace{0.1em}
  \begin{block}{默认路径的细节}
    \small
    \src{batch_build=True} 且传入 Python \src{op} 时，构造阶段主要使用 Torch 批量 Schur/stencil matvec；
    关闭批量模式或显式选择 C++ matvec 时，Level 0 使用 \src{CudaSchurOp}，后续层可使用
    \src{CudaCoarseSchurOp}。两者都实现同一个 $A_l$。
  \end{block}
  \source{\src{_multi_gpu.py:501,539}；\src{_schur_op.py:68,149}；\src{run_qcu_mg.py:164}。}
\end{frame}

\begin{frame}[t]{V-cycle 中算子的使用顺序与最终边界}
  \begin{columns}[T,totalwidth=\textwidth]
    \begin{column}{0.52\textwidth}
      \begin{tabular}{@{}l@{}}
        $r_l$：父层残差 \\
        $\quad\downarrow\quad R_l=P_l^\dagger$ \\
        $r_{l+1}$：粗层右端 \\
        $\quad\downarrow\quad A_{l+1}x_{l+1}=r_{l+1}$ \\
        $x_{l+1}$：递归粗层解 \\
        $\quad\downarrow\quad P_l$ \\
        $e_l$：父层修正 \\
        $\quad\downarrow\quad x_l\leftarrow x_l+e_l$，重算 $r_l$
      \end{tabular}
    \end{column}
    \begin{column}{0.44\textwidth}
      \begin{block}{层级关系}
        \small
        \begin{align*}
          A_0&=S,\\
          A_1&=P_0^\dagger S P_0,\\
          A_2&=P_1^\dagger A_1P_1.
        \end{align*}
        对任意更深层重复同样模式。
      \end{block}
      \begin{block}{边界}
        \footnotesize
        full-site 真残差与偶分量恢复只用于结果检查/输出，不参与粗算子构造；具体测例数值本次未重测。
      \end{block}
    \end{column}
  \end{columns}
  \source{\src{lattice_clover_multigrid.h:3348,4789,4846}。}
\end{frame}

\begin{frame}[t]{源码证据索引与范围说明}
  \begin{tabularx}{0.96\textwidth}{@{}P{0.27\linewidth}P{0.22\linewidth}Y@{}}
    \toprule
    证据对象 & 层次 & 关键定位 \\
    \midrule
    Level 0 Schur & CUDA-C++ 核心 & \src{lattice_clover_multigrid.h:815} \\
    层级参数 & CUDA-C++ 初始化 & \src{lattice_clover_multigrid.h:2643} \\
    null vector & Python 工具 & \src{_multigrid.py:355} \\
    局部 QR/transfer & Python/CUDA & \src{_multigrid.py:69}；\src{multigrid.cu:73} \\
    Galerkin probing & Python 工具 & \src{_multigrid.py:691} \\
    33 点 stencil & CUDA kernel & \src{multigrid.cu:858} \\
    指针接线 & C API bridge & \src{apply_clover_multigrid.cu:11} \\
    缓存/构造 & Python 驱动 & \src{_multi_gpu.py:490} \\
    \bottomrule
  \end{tabularx}
  \source{范围：相关入口、核心头文件、CUDA kernel、Python 构造器和缓存入口；源码关系与公式\verified{}，性能与具体 lattice 数值质量\unverified{}。}
\end{frame}

\end{document}
